% GENERATED FILE. Edit canonical lessons, then run npm run generate:books.
% canonical-source-hash: fnv1a64:2b1b160a8522af7a
% canonical-lessons: TA-C258-uruku, TA-C258-pari, TA-C258-kurai, TA-C258-terru, TA-C258-kai-kulukku

\chapter{To melt, To pick, To bark, To comfort, To shake hands}
\label{ch:ta-to-melt-to-pick-to-bark-to-comfort-to-shake-hands}
\hlchaptermodality{258}

\begin{chapteropening}
\textbf{By the end of this chapter:} \emph{I can say to melt, to pick, to bark, to comfort and to shake hands.}

Complete the last lesson of chapter 258: I can say to melt, to pick, to bark, to comfort and to shake hands.
\end{chapteropening}

\section[\texorpdfstring{\ta{உருகு} (uruku)}{உருகு (uruku)}]{\ta{உருகு} (uruku) — to melt}

\label{lesson:TA-C258-uruku}



\begin{quote}
Before the new one: say the Tamil for to be jealous, then the Tamil for to kiss.
\end{quote}

\subsection*{You'll want to know: \ta{உருகு}}
\textbf{\ta{உருகு}} — \emph{uruku} — \textquotedblleft{}to melt\textquotedblright{}.

\begin{grammarlens}[title={What you've built}]
One more word for this chapter.
\end{grammarlens}

\subsection*{Your turn}
\begin{itemize}
\raggedright
  \item \emph{Say it:} \emph{uruku}
  \item \emph{Say it:} \emph{uruku}, once more
  \item \emph{Say it:} say \emph{uruku}
  \item \emph{Recall:} say the Tamil for to be jealous, then the Tamil for to kiss, then say \emph{uruku} again
\end{itemize}

\begin{tcolorbox}[breakable,colback=teal!4,colframe=teal!35!black,title={Before you move on}]
What does \ta{உருகு} mean? (\textquotedblleft{}To melt\textquotedblright{}.) Say it once more.
\end{tcolorbox}
\section[\texorpdfstring{\ta{பறி} (paṟi)}{பறி (paṟi)}]{\ta{பறி} (paṟi) — to pick, to pluck}

\label{lesson:TA-C258-pari}



\begin{quote}
Before the new one: say the Tamil for to kiss, then the Tamil for to melt.
\end{quote}

\subsection*{You'll want to know: \ta{பறி}}
\textbf{\ta{பறி}} — \emph{paṟi} — \textquotedblleft{}to pick, to pluck\textquotedblright{}.

\begin{grammarlens}[title={What you've built}]
One more word for this chapter.
\end{grammarlens}

\subsection*{Your turn}
\begin{itemize}
\raggedright
  \item \emph{Say it:} \emph{paṟi}
  \item \emph{Say it:} \emph{paṟi}, once more
  \item \emph{Say it:} say \emph{paṟi}
  \item \emph{Recall:} say the Tamil for to kiss, then the Tamil for to melt, then say \emph{paṟi} again
\end{itemize}

\begin{tcolorbox}[breakable,colback=teal!4,colframe=teal!35!black,title={Before you move on}]
What does \ta{பறி} mean? (\textquotedblleft{}To pick, to pluck\textquotedblright{}.) Say it once more.
\end{tcolorbox}
\section[\texorpdfstring{\ta{குரை} (kurai)}{குரை (kurai)}]{\ta{குரை} (kurai) — to bark}

\label{lesson:TA-C258-kurai}



\begin{quote}
Before the new one: say the Tamil for to melt, then the Tamil for to pick.
\end{quote}

\subsection*{You'll want to know: \ta{குரை}}
\textbf{\ta{குரை}} — \emph{kurai} — \textquotedblleft{}to bark\textquotedblright{}.

\begin{grammarlens}[title={What you've built}]
One more word for this chapter.
\end{grammarlens}

\subsection*{Your turn}
\begin{itemize}
\raggedright
  \item \emph{Say it:} \emph{kurai}
  \item \emph{Say it:} \emph{kurai}, once more
  \item \emph{Say it:} say \emph{kurai}
  \item \emph{Recall:} say the Tamil for to melt, then the Tamil for to pick, then say \emph{kurai} again
\end{itemize}

\begin{tcolorbox}[breakable,colback=teal!4,colframe=teal!35!black,title={Before you move on}]
What does \ta{குரை} mean? (\textquotedblleft{}To bark\textquotedblright{}.) Say it once more.
\end{tcolorbox}
\section[\texorpdfstring{\ta{தேற்று} (tēṟṟu)}{தேற்று (tēṟṟu)}]{\ta{தேற்று} (tēṟṟu) — to comfort, to console}

\label{lesson:TA-C258-terru}



\begin{quote}
Before the new one: say the Tamil for to pick, then the Tamil for to bark.
\end{quote}

\subsection*{You'll want to know: \ta{தேற்று}}
\textbf{\ta{தேற்று}} — \emph{tēṟṟu} — \textquotedblleft{}to comfort, to console\textquotedblright{}.

\begin{grammarlens}[title={What you've built}]
One more word for this chapter.
\end{grammarlens}

\subsection*{Your turn}
\begin{itemize}
\raggedright
  \item \emph{Say it:} \emph{tēṟṟu}
  \item \emph{Say it:} \emph{tēṟṟu}, once more
  \item \emph{Say it:} say \emph{tēṟṟu}
  \item \emph{Recall:} say the Tamil for to pick, then the Tamil for to bark, then say \emph{tēṟṟu} again
\end{itemize}

\begin{tcolorbox}[breakable,colback=teal!4,colframe=teal!35!black,title={Before you move on}]
What does \ta{தேற்று} mean? (\textquotedblleft{}To comfort, to console\textquotedblright{}.) Say it once more.
\end{tcolorbox}
\section[\texorpdfstring{\ta{கை} \ta{குலுக்கு} (kai kulukku)}{கை குலுக்கு (kai kulukku)}]{\ta{கை} \ta{குலுக்கு} (kai kulukku) — to shake hands}

\label{lesson:TA-C258-kai-kulukku}



\begin{quote}
Before the new one: say the Tamil for to bark, then the Tamil for to comfort.
\end{quote}

\subsection*{You'll want to know: \ta{கை} \ta{குலுக்கு}}
\textbf{\ta{கை} \ta{குலுக்கு}} — \emph{kai kulukku} — \textquotedblleft{}to shake hands\textquotedblright{}.

\begin{grammarlens}[title={What you've built}]
One more word for this chapter.
\end{grammarlens}

\subsection*{Your turn}
\begin{itemize}
\raggedright
  \item \emph{Say it:} \emph{kai kulukku}
  \item \emph{Say it:} \emph{kai kulukku}, once more
  \item \emph{Say it:} say \emph{kai kulukku}
  \item \emph{Recall:} say the Tamil for to bark, then the Tamil for to comfort, then say \emph{kai kulukku} again
\end{itemize}

\begin{tcolorbox}[breakable,colback=teal!4,colframe=teal!35!black,title={Before you move on}]
What does \ta{கை} \ta{குலுக்கு} mean? (\textquotedblleft{}To shake hands\textquotedblright{}.) Say it once more.
\end{tcolorbox}
